% ============================================================
%  يُترجَم بـ XeLaTeX  (في Overleaf: Menu → Compiler → XeLaTeX)
%  الخط: Amiri (متوفر في Overleaf، أو حمّله من github.com/aliftype/amiri)
% ============================================================
\documentclass[12pt,a4paper]{article}
\usepackage[a4paper,margin=1.7cm,top=2.3cm,bottom=2.9cm,headheight=15pt,footskip=1.5cm]{geometry}
\usepackage{amsmath,amssymb}
\usepackage{icomma}
\usepackage[english,bidi=default,layout=lists]{babel}
\babelprovide[import,main]{arabic}
\babelfont{rm}[Scale=1.08]{Amiri}
\babelfont{sf}[Scale=1.08]{Amiri}
\usepackage[most]{tcolorbox}
\usepackage{tikz}
\usetikzlibrary{arrows.meta,calc,shadings}
\usepackage{array}
\usepackage{colortbl}
\usepackage{fancyhdr}
\usepackage{titlesec}
\usepackage{enumitem}
\usepackage{collcell}
\usepackage{needspace}
\usepackage{etoolbox}
% أعمدة جداول تُكتب خلاياها بالعربية في الاتجاه الصحيح
\newcolumntype{C}{>{\collectcell\ar}c<{\endcollectcell}}
% تجنب عنوان يتيم في أسفل الصفحة
\pretocmd{\section}{\Needspace{15\baselineskip}}{}{}
\pretocmd{\subsection}{\Needspace{11\baselineskip}}{}{}

% ---------- الألوان ----------
\definecolor{cDef}{HTML}{1F4E9C}
\definecolor{cThm}{HTML}{B3261E}
\definecolor{cRem}{HTML}{2E7D32}
\definecolor{cEx}{HTML}{C76B00}
\definecolor{cMeth}{HTML}{6A1B9A}
\definecolor{cSol}{HTML}{00695C}
\definecolor{cSec}{HTML}{1F3A68}

% ---------- العناوين ----------
\titleformat{\section}{\Large\bfseries\color{cSec}}{\thesection.}{0.6em}{}[{\color{cSec!60}\titlerule[1pt]}]
\titleformat{\subsection}{\large\bfseries\color{cThm!85!black}}{\thesubsection.}{0.5em}{}
\titlespacing*{\section}{0pt}{14pt}{8pt}
\titlespacing*{\subsection}{0pt}{10pt}{5pt}
\setlength{\parindent}{0pt}
\setlength{\parskip}{4pt}
\setlist{itemsep=2pt,topsep=3pt}

% ---------- الصناديق (تُرسم في سياق LTR ويُكتب محتواها بالعربية) ----------
\newtcolorbox{arbox}[3][]{enhanced,
  colback=#2!5,colframe=#2!80!black,arc=2mm,boxrule=0.8pt,
  left=6pt,right=6pt,top=8pt,bottom=5pt,
  fonttitle=\bfseries,coltitle=white,
  title={\foreignlanguage{arabic}{#3}},
  attach boxed title to top right={xshift=-5mm,yshift=-3mm},
  boxed title style={colback=#2!80!black,arc=1.5mm,boxrule=0pt},
  before upper={\selectlanguage{arabic}},
  before skip=8pt,after skip=8pt,breakable,#1}
\newenvironment{BOX}[2]{\begin{otherlanguage}{english}\begin{arbox}{#1}{#2}}{\end{arbox}\end{otherlanguage}}
\newenvironment{tarif}[1][تعريف]{\begin{BOX}{cDef}{#1}}{\end{BOX}}
\newenvironment{khassia}[1][خاصية]{\begin{BOX}{cThm}{#1}}{\end{BOX}}
\newenvironment{molahada}[1][ملاحظة]{\begin{BOX}{cRem}{#1}}{\end{BOX}}
\newenvironment{mithal}[1][مثال]{\begin{BOX}{cEx}{#1}}{\end{BOX}}
\newenvironment{tariqa}[1][طريقة]{\begin{BOX}{cMeth}{#1}}{\end{BOX}}

% ---------- التمارين والحلول ----------
\newcounter{exo}
\newenvironment{tamrin}[1]{\stepcounter{exo}\begin{BOX}{cSec}{التمرين \theexo: #1}}{\end{BOX}}
\newenvironment{hal}{\Needspace{12\baselineskip}\begin{otherlanguage}{english}\begin{arbox}[breakable]{cSol}{حل التمرين \theexo}}{\end{arbox}\end{otherlanguage}}

% رسومات TikZ تُوضع في سياق LTR
\newenvironment{LTR}{\begin{otherlanguage}{english}}{\end{otherlanguage}}
\newcommand{\ar}[1]{\foreignlanguage{arabic}{#1}}

% ---------- رموز ----------
% نقاط الحذف داخل الرياضيات (\dots يسبب خللاً في الاتجاه مع babel)
\renewcommand{\dots}{\mathinner{\ldotp\ldotp\ldotp}}
\renewcommand{\ldots}{\mathinner{\ldotp\ldotp\ldotp}}
\renewcommand{\cdots}{\mathinner{\cdotp\cdotp\cdotp}}
\newcommand{\N}{\mathbb{N}}
\newcommand{\Z}{\mathbb{Z}}
\newcommand{\D}{\mathbb{D}}
\newcommand{\Q}{\mathbb{Q}}
\newcommand{\R}{\mathbb{R}}
\newcommand{\pgcd}{\operatorname{PGCD}}
\newcommand{\ppcm}{\operatorname{PPCM}}

% ---------- تمثيل مجال على مستقيم عددي ----------
% \NLseg{a}{b}{[}{]}  مجال محدود
\newcommand{\NLseg}[4]{\begin{LTR}\begin{tikzpicture}[baseline=-0.6ex]
 \draw[-{Stealth[length=2mm]},gray!70](0,0)--(3.4,0);
 \draw[line width=2pt,cThm!80!black](0.7,0)--(2.7,0);
 \node[cThm!80!black,font=\LARGE] at (0.7,0){$#3$};
 \node[cThm!80!black,font=\LARGE] at (2.7,0){$#4$};
 \node[below=3pt,font=\footnotesize] at (0.7,0){$#1$};
 \node[below=3pt,font=\footnotesize] at (2.7,0){$#2$};
\end{tikzpicture}\end{LTR}}
% \NLright{a}{[}  مجال نحو +∞
\newcommand{\NLright}[2]{\begin{LTR}\begin{tikzpicture}[baseline=-0.6ex]
 \draw[-{Stealth[length=2mm]},gray!70](0,0)--(3.4,0);
 \draw[line width=2pt,cThm!80!black](1.3,0)--(3.3,0);
 \node[cThm!80!black,font=\LARGE] at (1.3,0){$#2$};
 \node[below=3pt,font=\footnotesize] at (1.3,0){$#1$};
\end{tikzpicture}\end{LTR}}
% \NLleft{b}{]}  مجال نحو −∞
\newcommand{\NLleft}[2]{\begin{LTR}\begin{tikzpicture}[baseline=-0.6ex]
 \draw[-{Stealth[length=2mm]},gray!70](0,0)--(3.4,0);
 \draw[line width=2pt,cThm!80!black](0.05,0)--(2.1,0);
 \node[cThm!80!black,font=\LARGE] at (2.1,0){$#2$};
 \node[below=3pt,font=\footnotesize] at (2.1,0){$#1$};
\end{tikzpicture}\end{LTR}}

% ---------- أيقونات مواقع التواصل (مرسومة بـ TikZ، لا تحتاج أي حزمة) ----------
\definecolor{cFB}{HTML}{1877F2}
\definecolor{cTG}{HTML}{229ED9}
\definecolor{cIG1}{HTML}{FEDA75}
\definecolor{cIG2}{HTML}{D62976}
\definecolor{cIG3}{HTML}{4F5BD5}
\newcommand{\fbicon}{\begin{LTR}\begin{tikzpicture}[baseline=-0.75ex,x=1em,y=1em]
 \fill[cFB] (0,0) circle (0.55);
 \node[white,inner sep=0,font=\usefont{TU}{lmss}{bx}{n}\small] at (0.04,-0.02) {f};
\end{tikzpicture}\end{LTR}}
\newcommand{\igicon}{\begin{LTR}\begin{tikzpicture}[baseline=-0.75ex,x=1em,y=1em]
 \shade[lower left=cIG1,lower right=cIG2,upper left=cIG2,upper right=cIG3,rounded corners=0.28em] (-0.55,-0.55) rectangle (0.55,0.55);
 \draw[white,line width=0.09em,rounded corners=0.18em] (-0.36,-0.36) rectangle (0.36,0.36);
 \draw[white,line width=0.09em] (0,0) circle (0.17);
 \fill[white] (0.22,0.22) circle (0.05);
\end{tikzpicture}\end{LTR}}
\newcommand{\tgicon}{\begin{LTR}\begin{tikzpicture}[baseline=-0.75ex,x=1em,y=1em]
 \fill[cTG] (0,0) circle (0.55);
 \fill[white] (-0.33,0.0) -- (0.32,0.27) -- (0.2,-0.3) -- (0.0,-0.15) -- (-0.08,-0.28) -- (-0.08,-0.08) -- cycle;
 \fill[cTG!70!black] (-0.08,-0.08) -- (0.25,0.2) -- (0.0,-0.15) -- (-0.08,-0.28) -- cycle;
\end{tikzpicture}\end{LTR}}
% سطر مواقع التواصل أسفل كل صفحة
% (يُكتب السطر بترتيب بصري من اليسار إلى اليمين: الأيقونة يمين كل اسم)
\newcommand{\PageBox}{\begin{LTR}\tikz[baseline=(p.base)]\node[fill=black,text=white,inner sep=2.5pt,minimum width=1.6em,minimum height=1.3em,font=\usefont{TU}{lmss}{bx}{n}\small](p){\ifnum\value{page}<10 0\fi\arabic{page}};\end{LTR}}
\newcommand{\SocialRow}{\begin{otherlanguage}{english}%
 \ar{\textbf{الأستاذ سوايسية هشام للرياضيات}}\ \tgicon\hfill
 \textbf{prof\_hicham\_math}\ \igicon\hfill
 \ar{\textbf{الأستاذ سوايسية هشام للرياضيات}}\ \fbicon\hspace{2.2em}\PageBox\end{otherlanguage}}
% ---------- رأس وتذييل الصفحة ----------
\pagestyle{fancy}
\fancyhf{}
\renewcommand{\headrulewidth}{0.6pt}
\renewcommand{\footrulewidth}{0.4pt}
% \SetHeader{يمين الرأس}{يسار الرأس}{يمين التذييل}{يسار التذييل}
\newcommand{\SetHeader}[2]{%
 \fancyhead[C]{\parbox{\textwidth}{\ar{#1\hfill #2}}}%
 \fancyfoot[C]{\parbox{\textwidth}{\small\SocialRow}}}

\usepackage{tkz-tab}
\usepackage{multicol}
% جدول تغيرات / إشارة (يُرسم في سياق LTR)
\newenvironment{TAB}{\par\begin{center}\begin{LTR}\begin{tikzpicture}[scale=0.95,every node/.style={font=\small}]}{\end{tikzpicture}\end{LTR}\end{center}}
% معلم بشبكة: \begin{GR}{xmin}{xmax}{ymin}{ymax}{unit cm}
\newenvironment{GR}[5]{\par\begin{center}\begin{LTR}\begin{tikzpicture}[x=#5cm,y=#5cm,font=\footnotesize]
 \def\xmn{#1}\def\xmx{#2}\def\ymn{#3}\def\ymx{#4}
 \draw[gray!25,very thin] (#1,#3) grid (#2,#4);
 \draw[-{Stealth[length=2mm]}] (#1,0)--(#2+0.3,0) node[below left]{$x$};
 \draw[-{Stealth[length=2mm]}] (0,#3)--(0,#4+0.3) node[below left]{$y$};
 \node[below left] at (0,0){$O$};
 \draw[very thick,-{Stealth[length=1.6mm]}] (0,0)--(1,0) node[below,pos=.55]{$\vec i$};
 \draw[very thick,-{Stealth[length=1.6mm]}] (0,0)--(0,1) node[left,pos=.55]{$\vec j$};
 \begin{scope}\clip (#1,#3) rectangle (#2,#4);}
{\end{scope}\end{tikzpicture}\end{LTR}\end{center}}
\newcommand{\Ticks}[2]{\foreach \t in {#1} \node[below,font=\scriptsize] at (\t,0){$\t$}; \foreach \t in {#2} \node[left,font=\scriptsize] at (0,\t){$\t$};}
\newcommand{\bx}[1]{\ensuremath{\colorbox{cDef!12}{$\displaystyle #1$}}}
\newcommand{\hl}[1]{\ensuremath{\colorbox{cEx!20}{$\displaystyle #1$}}}
\newcommand{\Lim}[1]{\displaystyle\lim_{x\to #1}}
\newcommand{\Cf}{(C_f)}
\newenvironment{qst}{\begin{enumerate}[label=\protect\circled{\arabic*},leftmargin=*]}{\end{enumerate}}
\newcommand{\circled}[1]{\foreignlanguage{english}{\tikz[baseline=(c.base)]\node[circle,fill=cDef,text=white,inner sep=1.2pt,font=\scriptsize\bfseries](c){#1};}}
% عمودان: \TwoCol{عرض اليسار}{محتوى اليسار}{عرض اليمين}{محتوى اليمين}  (آمن مع TikZ)
\newcommand{\TwoCol}[4]{\par\begin{LTR}\noindent\begin{minipage}[t]{#1\textwidth}\selectlanguage{arabic}#2\end{minipage}\hfill\begin{minipage}[t]{#3\textwidth}\selectlanguage{arabic}#4\end{minipage}\end{LTR}\par}
\newcommand{\un}{(u_n)}\newcommand{\vn}{(v_n)}
% برهان بالتراجع: \Rec{الخاصية P(n)}{التحقق من P(0)}{البرهان على الوراثة}
\newcommand{\Rec}[3]{\textbf{برهان بالتراجع:} نسمي $P(n)$ الخاصية: «#1».
\begin{itemize}
\item \textbf{التحقق:} #2
\item \textbf{الوراثة:} نفرض أن $P(n)$ صحيحة ونبرهن صحة $P(n+1)$: #3
\item \textbf{الخلاصة:} حسب مبدأ الاستدلال بالتراجع، $P(n)$ صحيحة من أجل كل عدد طبيعي $n$.
\end{itemize}}
\newcommand{\src}[1]{{\footnotesize\textbf{#1}}\par}
\newcommand{\ok}{$\textcolor{cRem}{\checkmark}$}
\newcommand{\limn}{\displaystyle\lim_{n\to+\infty}}
% أسئلة فرعية (أ) (ب) (جـ) كل سؤال في سطر
\newenvironment{sq}{\begin{itemize}[leftmargin=2.4em,itemsep=3pt]}{\end{itemize}}
\newcommand{\qa}{\item[(أ)]}\newcommand{\qb}{\item[(ب)]}\newcommand{\qc}{\item[(جـ)]}\newcommand{\qd}{\item[(د)]}\newcommand{\qe}{\item[(هـ)]}
\newcommand{\pa}{\par\smallskip\textbf{(أ)} }\newcommand{\pb}{\par\smallskip\textbf{(ب)} }\newcommand{\pc}{\par\smallskip\textbf{(جـ)} }\newcommand{\pd}{\par\smallskip\textbf{(د)} }\newcommand{\pe}{\par\smallskip\textbf{(هـ)} }
\newcommand{\ccl}[1]{\par\textbf{الخلاصة:} #1}

\pretocmd{\section}{\Needspace{24\baselineskip}}{}{}
\SetHeader{\small\bfseries\color{cSec}تمارين: المرجح في المستوي}{\small\color{cSec}السنة الثانية ثانوي -- الشعب العلمية}
\begin{document}
\newcolumntype{M}[1]{>{\centering\arraybackslash}p{#1}}
\newcommand{\RL}[4]{\TwoCol{#3}{#4}{#1}{#2}}
\newcommand{\QQ}[1]{\par\smallskip\textbf{#1.} }
\newcommand{\V}[1]{\overrightarrow{#1}}
\newcommand{\nV}[1]{\big\|#1\big\|}
\newcommand{\ABC}[3]{\\[2pt]\begin{tabular}{p{0.28\textwidth}p{0.28\textwidth}p{0.28\textwidth}}(جـ) \ #3 & (ب) \ #2 & (أ) \ #1\end{tabular}}
\begin{LTR}
\begin{tcolorbox}[enhanced,colback=cSec,colframe=cSec,arc=3mm,halign=center,boxrule=0pt,top=12pt,bottom=12pt,before upper={\selectlanguage{arabic}}]
{\color{white}\fontsize{26}{32}\selectfont\bfseries سلسلة تمارين: المُرَجِّح في المستوي}\\[8pt]
{\color{white}\large تمارين متنوعة ومسائل شاملة -- بدون حلول (للتدرّب)}\\[5pt]
{\color{white}\large السنة الثانية ثانوي -- الشعب العلمية}
\end{tcolorbox}
\end{LTR}

% =====================================================================
\section{مرجح نقطتين}
% ------------------------------------------------- 1
\begin{tamrin}{تعيين وإنشاء مرجح نقطتين}
$A$ و $B$ نقطتان متمايزتان من المستوي حيث $AB=6\ \text{cm}$.
\begin{enumerate}
\item أنشئ النقطة $G$ مرجح الجملة المثقلة $\{(A;2);(B;3)\}$.
\item أنشئ النقطة $N$ مرجح الجملة المثقلة $\{(A;-3);(B;1)\}$.
\item لتكن $K$ النقطة حيث $\V{BK}=-\dfrac12\V{AK}$. أثبت أن $K$ مرجح النقطتين $A$ و $B$ بمعاملين يطلب تعيينهما.
\item هل يوجد مرجح للجملة $\{(A;5);(B;-5)\}$؟ علّل.
\end{enumerate}
\end{tamrin}

% ------------------------------------------------- 2
\begin{tamrin}{علاقة شعاعية وتعيين المعاملات}
$A$ و $B$ نقطتان متمايزتان.
\begin{enumerate}
\item لتكن $G$ نقطة حيث $\V{AG}=\dfrac{3}{10}\V{AB}$. عيّن عددين حقيقيين $m_A$ و $m_B$ حتى تكون $G$ مرجح $\{(A;m_A);(B;m_B)\}$.
\item لتكن $I$ النقطة حيث $\V{AI}=\dfrac43\V{AB}$. بيّن أن $I$ مرجح النقطتين $A$ و $B$ ، ثم عيّن المعاملين.
\item بيّن أنه من أجل كل نقطة $M$: $\V{MA}+3\V{MB}=4\V{MJ}$ حيث $J$ نقطة يطلب تعيينها.
\end{enumerate}
\end{tamrin}

% =====================================================================
\section{مرجح ثلاث نقط وخاصية التجميع}
% ------------------------------------------------- 3
\begin{tamrin}{إنشاء مرجح ثلاث نقط بطريقتين}
$ABC$ مثلث ، و $G$ مرجح الجملة المثقلة $\{(A;2);(B;3);(C;1)\}$.
\begin{enumerate}
\item عبّر عن $\V{AG}$ بدلالة $\V{AB}$ و $\V{AC}$ ، ثم أنشئ النقطة $G$.
\item لتكن $I$ مرجح $\{(A;2);(B;3)\}$. بيّن أن $G$ مرجح $\{(I;5);(C;1)\}$ ، واستنتج طريقة ثانية لإنشاء $G$.
\item بيّن أنه من أجل كل نقطة $M$: \ $2\V{MA}+3\V{MB}-\V{MC}=4\V{MH}$ حيث $H$ نقطة يطلب تعيينها ، ثم استنتج أن $H$ مرجح $\{(I;5);(C;-1)\}$.
\end{enumerate}
\end{tamrin}

% ------------------------------------------------- 4
\begin{tamrin}{خاصية التجميع ومنتصف قطعة}
$ABC$ مثلث كيفي. $H$ نقطة تحقق العلاقة: $2\V{AH}+\V{AB}=\vec0$ ، و $G$ مرجح الجملة المثقلة $\{(A;3);(B;-1);(C;2)\}$.
\begin{enumerate}
\item بيّن أن $H$ هي مرجح النقطتين $A$ و $B$ المرفقتين على الترتيب بمعاملين يطلب تعيينهما.
\item عيّن ثم أنشئ النقطتين $H$ و $G$.
\item اعتماداً على خاصية التجميع أثبت أن $G$ منتصف القطعة المستقيمة $[HC]$.
\item $(E)$ مجموعة النقط $M$ من المستوي حيث: $\nV{3\V{MA}-\V{MB}+2\V{MC}}\le2\nV{\V{MA}-\V{MB}}$. \ عيّن ثم أنشئ المجموعة $(E)$.
\end{enumerate}
\end{tamrin}

% ------------------------------------------------- 5
\begin{tamrin}{تجميع بمعاملات سالبة}
$ABC$ مثلث ، و $G$ مرجح الجملة المثقلة $\{(A;-2);(B;1);(C;3)\}$.
\begin{enumerate}
\item أنشئ النقطة $G'$ مرجح الجملة المثقلة $\{(B;1);(C;3)\}$.
\item برهن أن $G$ مرجح الجملة المثقلة $\{(A;-2);(G';4)\}$.
\item أنشئ النقطة $G$.
\end{enumerate}
\end{tamrin}

% =====================================================================
\section{إحداثيا المرجح}
% ------------------------------------------------- 6
\begin{tamrin}{إحداثيات ومتوازي أضلاع}
في المستوي المنسوب إلى المعلم المتعامد والمتجانس $(O;\vec i,\vec j)$ ، نعتبر النقط: $A(1;0)$ ، $B(-3;2)$ و $C(1;3)$. \ ولتكن $G$ مرجح الجملة المثقلة $\{(A;-2);(B;1);(C;-1)\}$.
\begin{enumerate}
\item عيّن إحداثيي النقطة $G$.
\item عيّن إحداثيي النقطة $D$ حتى يكون الرباعي $ABCD$ متوازي أضلاع.
\item عيّن الأعداد الحقيقية $\alpha$ ، $\beta$ و $\gamma$ حتى تكون النقطة $D$ مرجح الجملة $\{(A;\alpha);(B;\beta);(C;\gamma)\}$.
\item عيّن $(\Gamma)$ مجموعة النقط $M$ من المستوي حيث: $\nV{-2\V{MA}+\V{MB}-\V{MC}}=2\nV{\V{MA}-\V{MB}+\V{MC}}$.
\end{enumerate}
\end{tamrin}

% ------------------------------------------------- 7
\begin{tamrin}{مرجح بوسيط ومحل هندسي}
في المستوي المنسوب إلى المعلم المتعامد والمتجانس $(O;\vec i,\vec j)$ ، نعتبر النقط $A(0;1)$ ، $B(0;5)$ و $C(2;1)$. \ لتكن $G_m$ مرجح الجملة المثقلة $\{(A;-m^2+4);(B;m^2-2m);(C;2m)\}$ حيث $m$ عدد حقيقي.
\begin{enumerate}
\item عيّن قيم $m$ التي من أجلها تكون $G_m$ موجودة ووحيدة.
\item عيّن إحداثيي النقطة $G_m$ بدلالة $m$.
\item ما هو المحل الهندسي للنقطة $G_m$ لما $m$ يمسح $\R$؟
\end{enumerate}
\end{tamrin}

% ------------------------------------------------- 8
\begin{tamrin}{استقامية ومجموعات نقط بالإحداثيات}
في المستوي المنسوب إلى المعلم المتعامد والمتجانس $(O;\vec i,\vec j)$ ، نعتبر النقط $A(1;2)$ ، $B(-8;-1)$ و $C(3;4)$ ، والنقطة $H$ المعرفة بالعلاقة $\V{AH}=\dfrac32\V{AC}$.
\begin{enumerate}
\item بيّن أن $H$ مرجح النقطتين $A$ و $C$ بمعاملين صحيحين يطلب تعيينهما ، ثم عيّن إحداثيي $H$.
\item عيّن إحداثيي النقطة $G$ مرجح الجملة المثقلة $\{(A;1);(B;-1);(C;-3)\}$.
\item بيّن أن النقط $B$ ، $H$ و $G$ على استقامة واحدة.
\item عيّن وأنشئ المجموعة $(E_1)$ للنقط $M$ حيث: $\nV{\V{MA}-\V{MB}-3\V{MC}}=6$.
\item عيّن المجموعة $(E_2)$ للنقط $M$ حيث: $2\nV{\V{MA}-\V{MB}-3\V{MC}}=3\nV{\V{MA}-3\V{MB}}$.
\end{enumerate}
\end{tamrin}

% ------------------------------------------------- 9
\begin{tamrin}{مرجح بوسيط $m$ ومركز ثقل}
المستوي منسوب إلى معلم متعامد ومتجانس $(O;\vec i,\vec j)$. نعتبر النقط $A(-2;0)$ ، $B(4;-2)$ و $C(2;3)$. \ $H$ نقطة من المستوي معرفة كما يلي: $-\V{HA}+2\V{BH}=\vec0$.\par
$G_m$ مرجح الجملة المثقلة $\{(A;m);(B;m^2+1);(C;4m-1)\}$.
\begin{enumerate}
\item بيّن أن النقطة $H$ هي مرجح للنقطتين $A$ و $B$ بمعاملين يطلب تعيينهما ، ثم احسب إحداثيي $H$.
\item ناقش حسب قيم العدد الحقيقي $m$ وجود النقطة $G_m$.
\item عيّن إحداثيي النقطة $G_1$ ، ثم بيّن أن $G_1$ هي مرجح للنقطتين $C$ و $H$. ماذا تستنتج بالنسبة للنقط $G_1$ ، $C$ و $H$؟
\item عيّن إحداثيي النقطة $I$ حتى تكون $C$ مركز ثقل المثلث $IG_2A$.
\item عيّن مجموعة النقط $M$ من المستوي التي تحقق:
\begin{sq}
\qa $\nV{\V{MA}+2\V{MB}+3\V{MC}}=3\nV{\V{MA}-\V{MB}}$.
\qb $\nV{\V{MA}+2\V{MB}+3\V{MC}}=2\nV{\V{MA}+2\V{MB}}$.
\end{sq}
\end{enumerate}
\end{tamrin}

% =====================================================================
\section{الاستقامية وتلاقي المستقيمات}
% ------------------------------------------------- 10
\begin{tamrin}{تلاقي ثلاثة مستقيمات}
$ABC$ مثلث ، ولتكن النقط $I$ ، $J$ و $K$ المعرفة كما يلي:
\begin{itemize}
\item $I$ مرجح الجملة المثقلة $\{(A;1);(B;-3)\}$.
\item $J$ تحقق العلاقة: $2\V{JA}-3\V{JC}=\vec0$.
\item $K$ تحقق العلاقة: $\V{BK}=\dfrac13\V{BC}$.
\end{itemize}
\begin{enumerate}
\item أنشئ شكلاً مناسباً.
\item أثبت أن $J$ و $K$ هما مرجحان لنقطتين من النقط $A$ ، $B$ و $C$ بمعاملين يطلب تعيينهما.
\item لتكن $G$ مرجح الجملة $\{(A;2);(B;-6);(C;-3)\}$. \ أثبت أن المستقيمات $(CI)$ ، $(BJ)$ و $(AK)$ تتلاقى في النقطة $G$.
\end{enumerate}
\end{tamrin}

% ------------------------------------------------- 11
\begin{tamrin}{استقامية ومحور قطعة}
$ABC$ مثلث ، $H$ نقطة معرفة بالعلاقة: $\V{HA}+3\V{HB}+2\V{BC}=\vec0$ ، و $G$ مرجح الجملة $\{(B;1);(C;2)\}$.
\begin{enumerate}
\item أنشئ النقطة $G$.
\item بيّن أن النقطة $H$ مرجح الجملة $\{(A;1);(B;1);(C;2)\}$ ، ثم أنشئها.
\item بيّن أن النقط $A$ ، $G$ و $H$ في استقامية.
\item $(\Gamma)$ مجموعة النقط $M$ من المستوي حيث: $\nV{\V{MB}+2\V{MC}}=\dfrac34\nV{\V{MA}+\V{MB}+2\V{MC}}$.
\begin{sq}
\qa عيّن وأنشئ المجموعة $(\Gamma)$.
\qb عيّن وأنشئ المجموعة $(\Gamma')$ حيث: $\nV{\V{MB}+2\V{MC}}\le\dfrac34\nV{\V{MA}+\V{MB}+2\V{MC}}$.
\end{sq}
\item $A(0;2)$ ، $B(-3;0)$ و $C(2;1)$ إحداثيات النقط في معلم متعامد ومتجانس.
\begin{sq}
\qa عيّن إحداثيي كل من النقطتين $G$ و $H$.
\qb تأكد حسابياً من صحة نتيجة السؤال 3.
\end{sq}
\end{enumerate}
\end{tamrin}

% ------------------------------------------------- 12
\begin{tamrin}{تقاطع مستقيمين ودائرة}
$ABC$ مثلث كيفي ، ولتكن $I$ نقطة تحقق: $-5\V{AI}+2\V{AC}=\vec0$ ، والنقطة $H$ مرجح الجملة المثقلة $\{(A;3);(B;-2)\}$.
\begin{enumerate}
\item بيّن أن $I$ مرجح النقطتين $A$ و $C$ بمعاملين يطلب تعيينهما.
\end{enumerate}
لتكن النقطة $G$ مرجح الجملة المثقلة $\{(A;3);(B;-2);(C;2)\}$.
\begin{enumerate}[resume]
\item أثبت أن $\V{IG}=-\dfrac23\V{IB}$. استنتج أن $G$ مرجح النقطتين $B$ و $I$ بمعاملين يطلب تعيينهما.
\item بيّن أن المستقيمين $(CH)$ و $(IB)$ متقاطعان في النقطة $G$.
\item عيّن $(C_1)$ مجموعة النقط $M$ من المستوي حيث: $\nV{3\V{MA}+2\V{MC}}=0$.
\item عيّن $(C_2)$ مجموعة النقط $M$ من المستوي حيث: $\nV{3\V{MA}-2\V{MB}+2\V{MC}}=\nV{\V{MC}-\V{MB}}$.
\end{enumerate}
\end{tamrin}

% ------------------------------------------------- 13
\begin{tamrin}{شعاع ثابت ومرجح بوسيط $k$}
\textbf{I)} $ABC$ مثلث من المستوي. $I$ و $J$ نقطتان حيث: $\V{AI}=\dfrac34\V{AB}$ و $6\V{BJ}-2\V{CJ}=\vec0$. \ $G$ مرجح الجملة المثقلة $\{(A;1);(B;3);(C;-1)\}$.
\begin{enumerate}
\item أنشئ النقطتين $I$ و $J$ ، ثم بيّن أن $I$ مرجح للنقطتين $A$ و $B$ يطلب تعيين معاملاتهما.
\item بيّن أن $G$ نقطة تلاقي المستقيمين $(CI)$ و $(AJ)$ ، ثم أنشئ النقطة $G$.
\item بيّن أن الشعاع $\vec v=\V{MA}+3\V{MB}-4\V{MC}$ مستقل عن $M$ ، ثم بيّن أن: $\vec v=4\V{CI}$.
\item عيّن المجموعة $(\Gamma_1)$ للنقط $M$ حيث: $\nV{\V{MA}+3\V{MB}}=\nV{6\V{MB}-2\V{MC}}$.
\item عيّن المجموعة $(\Gamma_2)$ للنقط $M$ حيث: $\nV{\V{MA}+3\V{MB}-\V{MC}}=\nV{\V{MA}+3\V{MB}-4\V{MC}}$.
\end{enumerate}
\textbf{II)} لتكن $G_k$ مرجح الجملة المثقلة $\{(A;k);(B;-k);(C;1)\}$ حيث $k$ عدد حقيقي.
\begin{enumerate}[resume]
\item هل النقطة $G_k$ موجودة من أجل كل عدد حقيقي $k$؟ علّل.
\item بيّن أن: $\V{CG_k}=-k\V{AB}$ ، ثم عيّن مجموعة النقط $G_k$ لما $k$ يمسح $\R$.
\item في المعلم المتعامد والمتجانس $(O;\vec i,\vec j)$ نعتبر $A(1;2)$ ، $B(0;-2)$ و $C(-1;3)$.
\begin{sq}
\qa عيّن بدلالة $k$ إحداثيي النقطة $G_k$.
\qb عيّن قيمة $k$ حتى تنتمي $G_k$ إلى المستقيم $(D)$ الذي معادلته $y=3x$.
\end{sq}
\end{enumerate}
\end{tamrin}

% ------------------------------------------------- 14
\begin{tamrin}{مجموعة نقط بتوظيف دراسة دالة}
$ABC$ مثلث و $m$ عدد حقيقي من المجال $[-1;1]$. \ نعتبر $G_m$ مرجح الجملة المثقلة $\{(A;m^2+1);(B;m);(C;-m)\}$.
\begin{enumerate}
\item برّر وجود $G_m$ من أجل كل عدد حقيقي $m$ من المجال $[-1;1]$ ، ثم بيّن أن: $\V{AG_m}=\dfrac{-m}{m^2+1}\V{BC}$.
\item أنشئ النقطتين $G_1$ و $G_{-1}$.
\item لتكن $N$ نقطة من $(BC)$. هل يمكن أن تنطبق $N$ على $G_m$؟ علّل.
\item شكّل جدول تغيرات الدالة $f$ المعرفة على $[-1;1]$ بـ: $f(x)=\dfrac{-x}{x^2+1}$.
\item استنتج مجموعة النقط $G_m$ عندما يمسح $m$ المجال $[-1;1]$.
\end{enumerate}
\end{tamrin}

% =====================================================================
\section{مجموعات النقط}
% ------------------------------------------------- 15
\begin{tamrin}{مثلث متقايس الأضلاع}
$ABC$ مثلث متقايس الأضلاع حيث $AB=AC=BC=2$. \ لتكن $(\Gamma)$ مجموعة النقط $M$ من المستوي التي تحقق:
\[\nV{\V{MA}-4\V{MB}+\V{MC}}=\nV{\V{MA}-2\V{MB}+\V{MC}}.\]
\begin{enumerate}
\item تحقق أن النقطة $B$ تنتمي إلى $(\Gamma)$.
\item بيّن أن الشعاع $\V{MA}-2\V{MB}+\V{MC}$ مستقل عن النقطة $M$.
\item ليكن $G$ مرجح الجملة المثقلة $\{(A;1);(B;-4);(C;1)\}$. بيّن أن $GM=\sqrt3$ من أجل كل $M\in(\Gamma)$.
\item استنتج طبيعة المجموعة $(\Gamma)$ محدداً عناصرها المميزة ، ثم أنشئها.
\end{enumerate}
\end{tamrin}

% ------------------------------------------------- 16
\begin{tamrin}{تطبيق فيزيائي: مركز الثقل (تمرين نموذجي)}
\begin{enumerate}
\item قضيب متجانس $[AB]$ كتلته مهملة وطوله $1{,}2\ \text{m}$. نعلّق كتلة $2\ \text{kg}$ في الطرف $A$ وكتلة $4\ \text{kg}$ في الطرف $B$. \ نعلم أن نقطة التوازن (مركز الثقل) $G$ هي مرجح $\{(A;2);(B;4)\}$. \ أين نعلّق القضيب حتى يبقى متوازناً؟
\item نضيف كتلة $3\ \text{kg}$ في منتصف القضيب $I$. عيّن الموضع الجديد لنقطة التوازن.
\item صفيحة رقيقة كتلتها مهملة على شكل مثلث $ABC$ ، حيث $A(0;0)$ ، $B(6;0)$ و $C(0;3)$ (الوحدة: dm). نضع كتلاً $1\ \text{kg}$ في $A$ ، $2\ \text{kg}$ في $B$ و $3\ \text{kg}$ في $C$. \ عيّن إحداثيي نقطة توازن الصفيحة.
\end{enumerate}
\end{tamrin}

% =====================================================================
\section{الاختيار من متعدد وصحيح أم خطأ}
% ------------------------------------------------- 17
\begin{tamrin}{اختيار من متعدد}
عيّن الاقتراح الصحيح الوحيد في كل حالة مع التعليل:
\begin{enumerate}
\item $G$ مرجح $\{(A;2);(B;-1)\}$. \ لدينا:
\ABC{$\V{AG}=-\V{AB}$}{$\V{AG}=\V{AB}$}{$\V{AG}=2\V{AB}$}
\item الجملة المثقلة $\{(A;3);(B;-3)\}$:
\ABC{مرجحها منتصف $[AB]$}{ليس لها مرجح}{مرجحها $A$}
\item $G$ مركز ثقل المثلث $ABC$. من أجل كل نقطة $M$: $\V{MA}+\V{MB}+\V{MC}$ يساوي:
\ABC{$\V{MG}$}{$3\V{MG}$}{$\vec0$}
\item $A(0;0)$ ، $B(3;0)$ و $C(0;6)$. مرجح $\{(A;1);(B;2);(C;3)\}$ هو:
\ABC{$(1;3)$}{$(3;1)$}{$(1;2)$}
\item مجموعة النقط $M$ حيث $\nV{\V{MA}+\V{MB}}=AB$ هي:
\ABC{دائرة قطرها $[AB]$}{محور $[AB]$}{دائرة مركزها $A$}
\item مجموعة النقط $M$ حيث $\nV{2\V{MA}-\V{MB}-\V{MC}}=4$ ($ABC$ مثلث) هي:
\ABC{دائرة}{المستوي كله أو المجموعة الخالية}{مستقيم}
\end{enumerate}
\end{tamrin}

% ------------------------------------------------- 18
\begin{tamrin}{صحيح أم خطأ}
أجب بصحيح أو خطأ مع التبرير:
\begin{enumerate}
\item إذا كانت $G$ مرجح $\{(A;\alpha);(B;\beta)\}$ حيث $\alpha>0$ و $\beta>0$ فإن $G\in[AB]$.
\item مرجح $\{(A;2);(B;6)\}$ هو نفسه مرجح $\{(A;1);(B;3)\}$.
\item الجملة المثقلة $\{(A;1);(B;1);(C;-2)\}$ تقبل مرجحاً.
\item إذا كانت $G$ مرجح $\{(A;1);(B;1);(C;2)\}$ و $I$ منتصف $[AB]$ فإن $G$ منتصف $[IC]$.
\item $\{(A;1);(B;-2)\}$ مرجحها $K$ و $I$ منتصف $[AB]$. مجموعة النقط $M$ حيث $\nV{\V{MA}-2\V{MB}}=\nV{\V{MA}+\V{MB}}$ هي محور $[KI]$.
\end{enumerate}
\end{tamrin}

\end{document}
